\section{Model}
\label{sec:model}

The economy consists of $S$ sectors, $R$ trade regions partitioned into
$B$ geopolitical blocs, a continuum of heterogeneous households indexed by
asset positions and labor types, a continuum of firms with pricing and
wage-setting power, and a monetary authority.

Time is discrete. The state of the economy at date $t$ includes the
distribution of household asset-type pairs $\Gamma_t(a, \ell)$, aggregate
productivity $\{z_{st}\}$, and the fragmentation premium $\phi_t$.

\subsection{Households}
\label{sec:households}

A continuum of households $i \in [0,1]$ are indexed by their labor type
$\ell \in \mathcal{L}$ (an occupation-location pair) and asset position
$a \in [\underline{a}, \infty)$. Households maximize:
\begin{equation}
  \max_{\{c_{it}, a_{it+1}, n_{it}\}} \mathbb{E}_0 \sum_{t=0}^{\infty} \beta^t
  \left[\frac{c_{it}^{1-\sigma}}{1-\sigma} - \chi \frac{n_{it}^{1+\varphi}}{1+\varphi}\right]
  \label{eq:hh_objective}
\end{equation}
subject to the budget constraint:
\begin{equation}
  c_{it} + a_{it+1} = (1 + r_t)\,a_{it} + w_t(\ell_i)\,n_{it} + T_{it}
  \label{eq:bc}
\end{equation}
and the natural borrowing constraint $a_{it+1} \geq \underline{a}$.

Here $r_t$ is the real interest rate, $w_t(\ell)$ is the effective wage for
labor type $\ell$ (reflecting the monopsony markdown, defined below), and
$T_{it}$ are lump-sum transfers. The distribution $\Gamma_t(a,\ell)$ is a
state variable that evolves endogenously.

\paragraph{Income risk.}
Labor income $w_t(\ell)\,n_{it}$ is subject to idiosyncratic income shocks.
We discretize the income process using the \citet{rouwenhorst1995} method
with $n_e = 7$ states, persistence $\rho_e = 0.966$, and innovation standard
deviation $\sigma_e = 0.5$ (in logs), following
\citet{floden2001idiosyncratic}.

\paragraph{Heterogeneous MPCs.}
A key property of the model is that the marginal propensity to consume (MPC)
varies across households. Define the household-specific MPC as:
\begin{equation}
  \mathcal{M}_i = \frac{\partial c_{it}}{\partial y_{it}}\Bigg|_{\text{transitory}}
\end{equation}
where $y_{it}$ is transitory income. Households near the borrowing constraint
(``hand-to-mouth'' households) have $\mathcal{M}_i \approx 1$; wealthy
households have $\mathcal{M}_i \approx 0$. The aggregate MPC is:
\begin{equation}
  \mathcal{M}_t = \int \mathcal{M}_i \, d\Gamma_t(a,\ell).
\end{equation}
This quantity responds to both monetary policy (via $r_t$) and the monopsony
wage markdown (via $w_t(\ell)$), a coupling absent from representative-agent
models and critical for the indirect transmission channel identified in
Proposition 3.

\subsection{Firms}
\label{sec:firms}

Firm $j$ in sector $s$, region $r$ produces output $y_{jt}$ using labor and
intermediate inputs.

\subsubsection{Production Technology}

Output is produced according to a nested CES technology:
\begin{equation}
  y_{jt} = z_{jt} \left[\alpha_n^{1/\theta} n_{jt}^{(\theta-1)/\theta} +
  \alpha_m^{1/\theta} M_{jt}^{(\theta-1)/\theta}\right]^{\theta/(\theta-1)}
  \label{eq:production}
\end{equation}
where $z_{jt}$ is idiosyncratic TFP, $n_{jt}$ is labor input, and
$M_{jt}$ is a CES aggregate of intermediate inputs from all sector-region
pairs:
\begin{equation}
  M_{jt} = \left[\sum_{s',r'} \omega_{s'r'}^{1/\eta}
  \left(\frac{m_{jt}^{s'r'}}{\tau_t^{rr'}}\right)^{(\eta-1)/\eta}
  \right]^{\eta/(\eta-1)}.
  \label{eq:composite_M}
\end{equation}
The term $\tau_t^{rr'} \geq 1$ is the iceberg trade cost between regions
$r$ and $r'$. The cost share parameters satisfy
$\alpha_n + \alpha_m = 1$. The elasticity $\theta$ governs substitution
between labor and intermediates; $\eta$ is the Armington elasticity across
source regions.

\subsubsection{Monopsony in Labor Markets}

Standard NK models assume firms are price-takers in the labor market. We
replace this assumption with an \emph{upward-sloping firm-level labor supply
curve}, which is the defining feature of oligopsonistic labor markets
\citep{manning2011, card2018}.

Firm $j$'s labor supply curve is:
\begin{equation}
  n_{jt} = \left(\frac{w_{jt}}{\bar{w}_t(\ell)}\right)^{\varepsilon_t}
  \bar{N}_t(\ell)
  \label{eq:labor_supply}
\end{equation}
where $\varepsilon_t > 0$ is the \emph{firm-level labor supply elasticity},
$\bar{w}_t(\ell)$ is the market-wide wage for labor type $\ell$, and
$\bar{N}_t(\ell)$ is aggregate labor supply of that type.

When $\varepsilon \to \infty$, equation \eqref{eq:labor_supply} delivers
perfect competition: the firm faces a perfectly elastic labor supply curve at
the market wage. When $\varepsilon$ is finite, the firm internalizes that
raising its wage attracts more workers, implying wage-setting power.

\paragraph{Countercyclical elasticity: search-friction microfoundation.}
\label{sec:micro_search}
Equation \eqref{eq:labor_supply} is not merely assumed---$\varepsilon_t$
is derived from search primitives following \citet{manning2003} and
\citet{berger2022monopsony}. In a directed-search model where employed workers
receive outside offers at rate $\lambda_e$, the firm's steady-state employment
equals inflows divided by outflows, and the resulting firm-level labor supply
elasticity is $\varepsilon_t = \varepsilon_\kappa[1 + \lambda_e(1-u_t)/u_t]$.
The Beveridge flow balance $\delta(1-u_t) = f(\theta_t)u_t$ implies the
offer-to-separation ratio falls when unemployment rises---reducing workers'
credible threat to quit and deepening monopsony power. A first-order Taylor
expansion yields:
\begin{equation}
  \varepsilon_t \approx \bar{\varepsilon} - \gamma\,(u_t - \bar{u}),
  \qquad
  \gamma \equiv \frac{\varepsilon_\kappa\,\lambda_e}{\bar{u}^2} > 0,
  \qquad \varepsilon_t \geq \varepsilon_{\min} = 0.5
  \label{eq:epsilon_cyclical}
\end{equation}
confirming $\gamma > 0$ by construction from search primitives, with
$\gamma$ calibrated to $0.8$ to match the observed labor share cyclicality
($-0.032$pp per 1pp rise in $u_t$, \citealt{rinz2022labor}).
The full derivation---including the structural formula
$\gamma = \varepsilon_\kappa\lambda_e/\bar u^2$ and its calibration to
matched employer-employee data---is in the Supplemental Appendix.

\paragraph{Wage-setting.}
Profit-maximizing wage-setting by firm $j$ yields the \emph{wage markdown
condition}:
\begin{equation}
  w_{jt}^* = \underbrace{\frac{\varepsilon_t}{\varepsilon_t + 1}}_{\equiv\,\mu_t^w} \cdot
  \mathrm{MPL}_{jt}
  \label{eq:wage_markdown}
\end{equation}
where $\mu_t^w \in (0,1)$ is the \emph{wage markdown}---the ratio of wage to
marginal product. Under perfect competition, $\mu_t^w = 1$. The parameter
$\gamma > 0$ ensures that $\mu_t^w$ is countercyclical: it falls (the markup
deepens) when unemployment rises.

\paragraph{Calibration of $\bar\varepsilon$.}
We set $\bar\varepsilon = 4$, implying a steady-state markdown of
$\mu^w = 4/5 = 0.80$---consistent with the central estimate of
\citet{manning2021monopsony} and the sector-level estimates of
\citet{dube2020monopsony}. This also pins down $\varepsilon_\kappa = \bar\varepsilon/(1+\bar\Psi) \approx 3.8$
following the microfoundation in Section~\ref{sec:micro_search}, given the calibrated $\lambda_e$.

\subsubsection{Price-Setting}

In goods markets, firms face a standard Calvo \citeyearpar{calvo1983} pricing
friction. A fraction $1-\xi$ of firms can optimally reset their price each
period. A firm that resets chooses $P_{jt}^*$ to maximize the present
discounted value of profits, yielding the standard condition:
\begin{equation}
  P_{jt}^* = \frac{\varepsilon_p}{\varepsilon_p - 1} \cdot
  \mathbb{E}_t\sum_{k=0}^{\infty}(\xi\beta)^k\,\frac{\Lambda_{t+k}}{\Lambda_t}
  \cdot \widehat{mc}_{jt+k} \cdot y_{jt+k|t}
  \label{eq:optimal_price}
\end{equation}
where $\varepsilon_p$ is the elasticity of substitution across goods varieties,
$\Lambda_t$ is the stochastic discount factor, and $\widehat{mc}_{jt}$ is
the \emph{true} marginal cost under monopsonistic labor hiring:
\begin{equation}
  \widehat{mc}_{jt} = \frac{w_{jt}}{\mu_t^w \cdot z_{jt} \cdot f_n}
  = \frac{\mathrm{MPL}_{jt}}{z_{jt} \cdot f_n}.
  \label{eq:true_mc}
\end{equation}
Note that the markdown $\mu_t^w$ cancels in \eqref{eq:true_mc}: the true
marginal cost of an additional unit of output equals the marginal product
value regardless of wage-setting power. However, the markdown \emph{does}
affect real marginal cost dynamics through the general equilibrium link between
employment, productivity, and the markup.

\subsection{The Multi-Region IO Network}
\label{sec:io_network}

\subsubsection{Geopolitical Blocs and Trade Costs}

The world is partitioned into $B = 2$ geopolitical blocs:
$\mathcal{B}_1 = \{\text{US}, \text{EU}, \text{RoW}\}$ and
$\mathcal{B}_2 = \{\text{China}, \text{Emerging Asia}\}$. The iceberg trade
cost between regions $r$ and $r'$ is:
\begin{equation}
  \tau_t^{rr'} = \begin{cases}
    1 & \text{if } r = r' \\
    1 + \delta^{\text{within}} & \text{if } b(r) = b(r') \neq r \\
    1 + \delta^{\text{within}} + \underbrace{\phi_t}_{\text{fragmentation premium}}
      & \text{if } b(r) \neq b(r')
  \end{cases}
  \label{eq:trade_costs}
\end{equation}
The fragmentation premium $\phi_t$ follows an AR(1) process:
\begin{equation}
  \phi_t = \rho_\phi \phi_{t-1} + \sigma_\phi\,\varepsilon_t^\phi,
  \qquad \varepsilon_t^\phi \sim \mathcal{N}(0,1)
  \label{eq:phi_process}
\end{equation}
with high persistence $\rho_\phi = 0.90$, reflecting the slow-moving nature
of geopolitical realignments. The Supplemental Appendix provides a
strategic-game micro-foundation and demonstrates that the AR(1) is the
Nash-equilibrium linearization of a two-bloc tariff game.

\paragraph{Sector-specific regulatory premium.}
The empirical analysis (Section~\ref{sec:ip_frag}) documents that modern
geoeconomic fragmentation operates through IP and regulatory channels in
addition to tariff costs. We extend \eqref{eq:trade_costs} to accommodate
a sector-specific regulatory premium $\psi_{s,t}^{rr'}$:
\begin{equation}
  \tau_{s,t}^{rr'} = \begin{cases}
    1 & r = r' \\
    1 + \delta^{\text{within}} & b(r)=b(r'), r\neq r' \\
    1 + \delta^{\text{within}} + \phi_t + \psi_{s,t}^{rr'} & b(r)\neq b(r')
  \end{cases}
  \label{eq:trade_costs_full}
\end{equation}
In the baseline model, $\psi_{s,t} \equiv 0$ (sector-neutral $\phi_t$ shock).
In the extended model with IP decoupling (Supplemental Appendix
and Section~\ref{sec:ip_frag}), $\psi_{s,t}^{\text{Tech}}$ captures
semiconductor export control shocks. All main results hold under either
specification; the IP extension is used in the broadened fragmentation
exposure index in Section~\ref{sec:data}.

\subsubsection{The Leontief IO Structure}

We represent the global production network as a weighted directed graph
$\mathcal{G} = (\mathcal{V}, \mathcal{E}, \boldsymbol{\Omega})$ where:
\begin{itemize}
  \item Nodes $\mathcal{V}$: sector-region pairs $(s,r)$, with
    $|\mathcal{V}| = S \times R = 10 \times 5 = 50$.
  \item Edges $\mathcal{E}$: input-output linkages between node pairs.
  \item Weights $\Omega_{(s,r)(s',r')} \in [0,1]$: cost share of sector
    $(s',r')$ inputs in the production of sector $(s,r)$.
\end{itemize}

The IO balance system in matrix form is:
\begin{equation}
  \mathbf{X} = \boldsymbol{\Omega}^\top \mathbf{X} + \mathbf{F}
  \label{eq:io_balance}
\end{equation}
where $\mathbf{X} \in \mathbb{R}^{S \times R}$ is the gross output vector and
$\mathbf{F}$ is final demand. Solving \eqref{eq:io_balance}:
\begin{equation}
  \mathbf{X} = (\mathbf{I} - \boldsymbol{\Omega}^\top)^{-1}\mathbf{F}
  \equiv \mathbf{L}\,\mathbf{F}
  \label{eq:leontief}
\end{equation}
where $\mathbf{L}$ is the \emph{Leontief inverse}. Its $(i,j)$ entry gives
the total output of sector $i$ required per unit of final demand in sector $j$,
including all upstream supply chain rounds.

\paragraph{Domar Weights.}
Following \citet{baqaee2019}, the \emph{Domar weight} of sector $(s,r)$ is:
\begin{equation}
  \lambda_{sr} = \frac{P_{sr} X_{sr}}{P \cdot \mathrm{GDP}}
  \label{eq:domar}
\end{equation}
which equals the sector's share of \emph{gross} (not value-added) output in
GDP. Because sectors supply both final goods and intermediates, Domar weights
sum to more than 1: $\sum_{sr} \lambda_{sr} > 1$.

\paragraph{Fragmentation and Domar Weights.}
A key result of the model is that fragmentation reshapes the Domar weight
distribution. Specifically:

\begin{lemma}[Domar Weight Reallocation]
  \label{lemma:domar}
  For any upstream sector $(s,r)$ that primarily supplies cross-bloc
  intermediate goods: $\partial\lambda_{sr}/\partial\phi > 0$.
  For downstream service sectors with limited cross-bloc linkages:
  $\partial\lambda_{sr}/\partial\phi \approx 0$.
\end{lemma}

\begin{proof}
  See Appendix \ref{app:proofs}.
\end{proof}

The intuition is straightforward: when cross-bloc trade costs rise, domestic
firms substitute toward domestic upstream suppliers. This increases the volume
of domestic upstream production relative to GDP, raising upstream Domar
weights. The effect is amplified by the Leontief multiplier, which propagates
first-order substitutions into higher-order input demand changes.

\subsection{Monetary Authority}
\label{sec:monetary}

The central bank sets the nominal interest rate according to a (potentially
augmented) Taylor rule:
\begin{equation}
  i_t = r^* + \pi^* + \phi_\pi(\pi_t - \pi^*) + \phi_y\,\tilde{y}_t
        + \phi_\mu\,\hat\mu_t^w + \phi_\phi\,\hat\phi_t + \nu_t
  \label{eq:taylor}
\end{equation}
where $\tilde{y}_t$ is the output gap, $\hat\mu_t^w \equiv \mu_t^w - \bar\mu^w$
is the wage markdown gap, $\hat\phi_t \equiv \phi_t - \bar\phi$ is the
fragmentation premium deviation from steady state, and $\nu_t$ is an
exogenous monetary policy shock. The \emph{standard} Taylor rule corresponds
to $\phi_\mu = \phi_\phi = 0$; the \emph{augmented} rule allows these
coefficients to differ from zero.

Proposition 3 characterizes the welfare-optimal values of all four
coefficients $(\phi_\pi^*, \phi_y^*, \phi_\mu^*, \phi_\phi^*)$.
